\section{Circuit Complexity --- Karp Lipton Prefix}
%%% generated by the blueprint pipeline from TCSlib.Complexity.CircuitComplexity.KarpLiptonPrefix
\subsection*{Circuit Complexity --- Karp Lipton Prefix: Introduction}

Declarations of the Lean module \texttt{TCSlib.Complexity.CircuitComplexity.KarpLiptonPrefix}.

\subsection*{Circuit Complexity --- Karp Lipton Prefix: Results}

\begin{definition}[Truncating Boolean operation]
\lean{Complexity.KarpLipton.dropTr}
\label{Complexity.KarpLipton.dropTr}
\leanok
Define a Boolean operation of two inputs by returning the negation of the second input
when the first is false, and returning true when the first is true.
\end{definition}

\begin{definition}[Dropout transition rule]
\lean{Complexity.KarpLipton.dropOut}
\label{Complexity.KarpLipton.dropOut}
\leanok
The dropout transition takes a pair of Boolean values and returns no output when the
first value is false; when the first value is true, it returns the second value as its
output.
\end{definition}

\begin{definition}[Drop-marker transduction]
\lean{Complexity.KarpLipton.dropMarker}
\label{Complexity.KarpLipton.dropMarker}
\leanok
\uses{Complexity.KarpLipton.dropOut, Complexity.KarpLipton.dropTr, Complexity.transduce}
For a Boolean list $w$, the drop-marker operation is the list produced by applying the
transducer with transition rule that emits nothing on a false input and emits the input
value on a true input, starting from the state false. The transition function used by
the transducer is not specified in the provided statement.
\end{definition}

\begin{theorem}[Identity action from the true state]
\lean{Complexity.KarpLipton.transduce_drop_true}
\label{Complexity.KarpLipton.transduce_drop_true}
\leanok
\uses{Complexity.KarpLipton.dropOut, Complexity.KarpLipton.dropTr, Complexity.transduce}
\difficulty{4}
For every finite binary string $p$, the one-pass transducer with transition function
$\delta(q,b)=\neg b$ when $q$ is false and $\delta(q,b)=\mathrm{true}$ when $q$ is true,
and output function that emits nothing when $q$ is false and emits $b$ when $q$ is true,
started in the true state, outputs $p$ unchanged.
\end{theorem}

\begin{theorem}[Removing a marker prefix]
\lean{Complexity.KarpLipton.dropMarker_marker}
\label{Complexity.KarpLipton.dropMarker_marker}
\leanok
\uses{Complexity.KarpLipton.dropMarker, Complexity.KarpLipton.dropOut, Complexity.KarpLipton.dropTr, Complexity.KarpLipton.transduce_drop_true, Complexity.transduce}
\difficulty{4}
For every natural number $k$ and bit string $p$, applying the marker-deletion transducer
to the string consisting of $k$ ones, followed by $0$ and then $p$, returns $p$.
\end{theorem}

\begin{theorem}[Polynomial-Time Computability of Marker Dropping]
\lean{Complexity.KarpLipton.polyTimeComputable_dropMarker}
\label{Complexity.KarpLipton.polyTimeComputable_dropMarker}
\leanok
\uses{Complexity.KarpLipton.dropMarker, Complexity.KarpLipton.dropOut, Complexity.KarpLipton.dropTr, Complexity.PolyTimeComputable, Complexity.polyTimeComputable_transduce}
\difficulty{1}
The function on bit strings that processes each input bit with a two-state transducer,
initially in state $0$, is polynomial-time computable: in state $0$ it emits nothing and
moves to state $1$, while in state $1$ it emits the current bit and remains in state
$1$.
\end{theorem}

\begin{definition}[Witness length bound]
\lean{Complexity.KarpLipton.wLen}
\label{Complexity.KarpLipton.wLen}
\leanok
\uses{Complexity.pairFstD}
For natural numbers $C,c$ and a Boolean list $y$, define
$w(C,c,y)=C\bigl(\operatorname{length}(\operatorname{pairFstD}(y))+1\bigr)^c$, where
$\operatorname{pairFstD}(y)$ denotes the first component decoded from $y$.
\end{definition}

\begin{theorem}[Length from the first encoded component]
\lean{Complexity.KarpLipton.wLen_pairEncode}
\label{Complexity.KarpLipton.wLen_pairEncode}
\leanok
\uses{Complexity.KarpLipton.wLen, Turing.pairEncode}
\difficulty{1}
For all natural numbers $C,c$ and binary strings $x,u$, the length function applied to
the pair encoding of $x$ and $u$ is $C(|x|+1)^c$.
\end{theorem}

\begin{definition}[Prefix language from a length function]
\lean{Complexity.KarpLipton.prefixLangOf}
\label{Complexity.KarpLipton.prefixLangOf}
\leanok
\uses{Complexity.KarpLipton.dropMarker, Complexity.pairFstD, Complexity.pairSndD, Turing.pairEncode}
Given a function $\ell$ from Boolean lists to natural numbers and a language $V$ of
Boolean lists, a Boolean list $q$ belongs to the resulting language exactly when there
is a Boolean list $s$ such that appending $s$ to the dropout-transduced second component
decoded from $q$ gives a list of length $\ell$ of the first component decoded from $q$,
and encoding that first component together with the appended list produces a member of
$V$.
\end{definition}

\begin{theorem}[Membership in a Prefix Language]
\lean{Complexity.KarpLipton.mem_prefixLangOf}
\label{Complexity.KarpLipton.mem_prefixLangOf}
\leanok
\uses{Complexity.KarpLipton.dropMarker, Complexity.KarpLipton.prefixLangOf, Complexity.pairFstD, Complexity.pairSndD, Turing.pairEncode}
\difficulty{1}
Let $\ell$ assign a natural number to each bit string, let $V$ be a language of bit
strings, and let $q$ be a bit string. Then $q$ belongs to the language consisting of
strings whose decoded first component determines the length bound and whose decoded
second component can be extended after marker-dropping into a word of $V$ exactly when
there is a bit string $s$ such that $\lvert d(\pi_2(q))s\rvert=\ell(\pi_1(q))$ and
$\langle\pi_1(q),d(\pi_2(q))s\rangle\in V$, where $\pi_1,\pi_2$ are the total pair
projections and $d$ is the marker-dropping transduction.
\end{theorem}

\begin{theorem}[Prefix Membership After a Marker Run]
\lean{Complexity.KarpLipton.mem_prefixLangOf_iff}
\label{Complexity.KarpLipton.mem_prefixLangOf_iff}
\leanok
\uses{Complexity.KarpLipton.dropMarker_marker, Complexity.KarpLipton.mem_prefixLangOf, Complexity.KarpLipton.prefixLangOf, Complexity.pairFstD_pairEncode, Complexity.pairSndD_pairEncode, Turing.pairEncode}
\difficulty{3}
Let $\ell$ be a function from bit strings to natural numbers, let $V$ be a language of
bit strings, let $y,p$ be bit strings, and let $k$ be a natural number. The encoding of
$y$ paired with a string of $k$ ones followed by $0p$ belongs to the prefix language
defined by $\ell$ and $V$ if and only if there exists a bit string $s$ such that $\lvert
ps\rvert=\ell(y)$ and $\langle y,ps\rangle\in V$.
\end{theorem}

\begin{theorem}[Polynomially bounded prefix language in NP]
\lean{Complexity.KarpLipton.prefixLangOf_mem_NP}
\label{Complexity.KarpLipton.prefixLangOf_mem_NP}
\leanok
\uses{Complexity.KarpLipton.dropMarker, Complexity.KarpLipton.mem_prefixLangOf, Complexity.KarpLipton.polyTimeComputable_dropMarker, Complexity.KarpLipton.prefixLangOf, Complexity.NP, Complexity.P, Complexity.PolyHierarchy.UnaryPT, Complexity.PolyHierarchy.UnaryPT.comp, Complexity.PolyTimeComputable, Complexity.PolyTimeComputable.append, Complexity.PolyTimeComputable.comp, Complexity.PolyTimeComputable.pairEncode, Complexity.inter_mem_P, Complexity.lenEq_preimage_mem_P, Complexity.length_pairFstD_le, Complexity.mem_NP_iff_exists_length_le, Complexity.pairFstD, Complexity.pairFstD_pairEncode, Complexity.pairSndD, Complexity.pairSndD_pairEncode, Complexity.polyTimeComputable_pairFstD, Complexity.polyTimeComputable_pairSndD, Complexity.preimage_mem_P, Turing.pairEncode}
\difficulty{5}
Let $\ell$ be a function from binary strings to natural numbers, and let $V$ be a
language in $\mathsf{P}$. Suppose the unary language consisting of the strings of
$\ell(y)$ ones is polynomial-time computable, and suppose there are natural numbers
$A,a$ such that $\ell(y)\le A(|y|+1)^a$ for every binary string $y$. Then the language
of strings $q$ for which there is a binary string $s$ such that the marker-removed
second component of $q$, followed by $s$, has length $\ell$ of the first component of
$q$, and the pair encoding of that first component with this concatenation belongs to
$V$, is in $\mathsf{NP}$.
\end{theorem}

\begin{definition}[Prefix language construction]
\lean{Complexity.KarpLipton.prefixLang}
\label{Complexity.KarpLipton.prefixLang}
\leanok
\uses{Complexity.KarpLipton.prefixLangOf, Complexity.KarpLipton.wLen}
For natural numbers $C$ and $c$ and a language $V$ of Boolean strings, define the prefix
language of $V$ using the length bound determined by $C$ and $c$.
\end{definition}

\begin{theorem}[Membership in the Prefix Language]
\lean{Complexity.KarpLipton.mem_prefixLang}
\label{Complexity.KarpLipton.mem_prefixLang}
\leanok
\uses{Complexity.KarpLipton.dropMarker, Complexity.KarpLipton.prefixLang, Complexity.KarpLipton.wLen, Complexity.pairFstD, Complexity.pairSndD, Turing.pairEncode}
\difficulty{1}
For natural numbers $C,c$, a language $V$ of bit strings, and a bit string $q$,
membership of $q$ in the prefix language with length bound $C(\lvert\pi_1(q)\rvert+1)^c$
is equivalent to the existence of a bit string $s$ such that $\lvert
d(\pi_2(q))s\rvert=C(\lvert\pi_1(q)\rvert+1)^c$ and the encoded pair
$\langle\pi_1(q),d(\pi_2(q))s\rangle$ belongs to $V$. Here $\pi_1$ and $\pi_2$ are the
total pair projections, defaulting to the empty string when parsing fails, and $d$ is
the marker-dropping transduction.
\end{theorem}

\begin{theorem}[Polynomial Prefix Membership Criterion]
\lean{Complexity.KarpLipton.mem_prefixLang_iff}
\label{Complexity.KarpLipton.mem_prefixLang_iff}
\leanok
\uses{Complexity.KarpLipton.mem_prefixLangOf_iff, Complexity.KarpLipton.prefixLang, Complexity.KarpLipton.wLen, Turing.pairEncode}
\difficulty{1}
Let $C,c$ be natural numbers, let $V$ be a language of bit strings, let $y,p$ be bit
strings, and let $k$ be a natural number. The encoding of $y$ paired with a string of
$k$ ones followed by $0p$ belongs to the prefix language with length bound
$C(\lvert\pi_1(y)\rvert+1)^c$ if and only if there exists a bit string $s$ such that
$\lvert ps\rvert=C(\lvert\pi_1(y)\rvert+1)^c$ and $\langle y,ps\rangle\in V$, where
$\pi_1$ is the total first pair projection.
\end{theorem}

\begin{theorem}[Polynomial-length prefix language in NP]
\lean{Complexity.KarpLipton.prefixLang_mem_NP}
\label{Complexity.KarpLipton.prefixLang_mem_NP}
\leanok
\uses{Complexity.KarpLipton.prefixLang, Complexity.KarpLipton.prefixLangOf_mem_NP, Complexity.NP, Complexity.P, Complexity.PolyHierarchy.unaryPT_poly, Complexity.length_pairFstD_le, Complexity.polyTimeComputable_pairFstD}
\difficulty{3}
Let $C,c$ be natural numbers and let $V$ be a language in $\mathsf{P}$. Then the
language of binary strings $q$ for which there is a binary string $s$ such that the
marker-removed second component of $q$, followed by $s$, has length $C(|x|+1)^c$, where
$x$ is the first component of $q$, and the pair encoding of $x$ with this concatenation
belongs to $V$, is in $\mathsf{NP}$.
\end{theorem}
